\documentclass[10pt,a4paper]{amsart}
\usepackage[margin=19mm]{geometry}
\usepackage{amsmath,amssymb,mathtools}
\usepackage[scr=rsfs,cal=euler]{mathalfa}
\usepackage{xfrac}
\usepackage[unicode,hidelinks,bookmarks=false]{hyperref}
\newcommand{\ie}{i.e.\ }
\newcommand{\cf}{cf.\ }
\newcommand{\resp}{respectively\ }
\newcommand{\Subsection}{\subsection}
\providecommand{\nobreakdash}{\nobreak\hbox{-}}
\setlength{\emergencystretch}{2em}
\multlinegap0pt
\usepackage{bm}
\usepackage{bbm}
\usepackage{dsfont}
\usepackage{xspace}
\usepackage{cancel}
\usepackage{mathtools}
\usepackage{amsthm}
\makeatletter
\@ifundefined{assumption}{\newtheorem{assumption}{Assumption}}{}
\@ifundefined{proposition}{\newtheorem{proposition}{Proposition}}{}
\makeatother
\title[Convergence analysis for a tree-based nonlinear reduced basi]{Convergence analysis for a tree-based nonlinear reduced basis method}
\author{Mohamed Barakat, Diane Guignard}
\date{Source: arXiv:2511.00226v1 (CC BY 4.0)}
\begin{document}
\maketitle

\noindent\emph{Source and licence.} Convergence analysis for a tree-based nonlinear reduced basis method. Version arXiv:2511.00226v1 (CC BY 4.0), distributed under the Creative Commons Attribution 4.0 International licence, as recorded in the arXiv licence metadata for this exact version and shown on the abstract page https://arxiv.org/abs/2511.00226v1. Full rights notice: licenses/arxiv-2511.00226-CC-BY-4.0.txt.\par\smallskip

\noindent\textit{Editorial scope.} Counted unit: the Proposition whose source label is \texttt{hat C} in the article (source0.tex lines 447--453, source label \texttt{hat C}), delivered together with the article's own complete original proof of it (source0.tex lines 455--535, one \texttt{proof} environment of 81 lines). No mathematics is written, repaired or re-derived: statement and proof are the article's own text line for line. Required context retained uncounted: para pde (lines 68--71); assump para (lines 336--342); h (lines 435--438). Every definition, display and label of those lines is retained unchanged. Omitted from this package and not redistributed: 1--67, 72--335, 343--434, 439--446, 454--454, 536--986. Nothing inside a retained excerpt is omitted or altered: every retained body is delivered exactly as the article's lines are written, so no omission receipt of any kind accompanies this package. Citations: the retained excerpts carry no citation command at all, so no bibliography entry of the article is transcribed and no reference is redistributed. Reference adaptation: none was needed and none was made; every reference inside the delivered unit resolves inside it, so no unresolved reference or citation is produced. Printed slips kept literal: no printed word, symbol, hypothesis or number of the counted statement or of its proof was altered. Layout and class: the article's own class is \texttt{article} with packages including graphicx, amssymb, graphics, epsfig, float, subcaption, amsmath, amsthm, fontenc, accents, babel, xcolor, framed, soul; the wrapper is the lane's pinned \texttt{amsart} editorial wrapper at 10pt on A4, which restates the theorem-like environments the retained text needs and adds the article's own macro definitions that the retained text needs (none), with extra wrapper packages (bm, bbm, dsfont, xspace, cancel, mathtools). The delivered numbering is the unit's own. Assets: no retained span contains a figure environment, a graphics-inclusion command, a PDF-page inclusion, a table or a TikZ picture; no image file is copied into this package, the package declares no asset dependency and no image file is redistributed with it. Licence: version arXiv:2511.00226v1 of this article is distributed under the Creative Commons Attribution 4.0 International licence, as recorded in the arXiv licence metadata for this exact version and shown on the abstract page https://arxiv.org/abs/2511.00226v1. A full rights notice is delivered at licenses/arxiv-2511.00226-CC-BY-4.0.txt. Characters: every retained excerpt is pure ASCII, so the delivered engine, XeLaTeX with its default Latin Modern text font, reports no missing character.
\begin{equation}
    a(u(\mu), v ; \mu)=f(v ; \mu), \quad \forall v \in V.
    \label{para pde}
\end{equation}

\begin{assumption}
The parameter-to-solution map $\mu \mapsto u_{\mathcal{N}}(\mu)$ admits an extension to an open set $\mathcal{O} \subset \mathbb{C}^d$ containing the parameter domain $\mathcal{D}$ such that, for any $z \in \mathcal{O}$, the map $z \mapsto u_{\mathcal{N}}(z)$ is holomorphic in each variable $z_j$ with the uniform bound
\begin{equation}
    \sup_{z \in \mathcal{O}} \|u_{\mathcal{N}}(z)\|_V \leq C.
\end{equation}
\label{assump para}
\end{assumption}

\begin{equation}
    h:=\max_{\substack{1 \leq k \leq K \\ 1 \leq j \leq d}}h_{k,j}
    \label{h}
\end{equation}

\begin{proposition} 
    Consider an elliptic parametric PDE of the form \eqref{para pde} with a $d$-dimensional parameter domain such that the parameter-to-solution map satisfies Assumption \ref{assump para}. Furthermore, let $\mathcal{D}$ be partitioned into $K$ tensor-product-structured subdomains ${\mathcal{D}_k}$,$1 \leq k \leq K$, and let $h$ denote the maximum side length of the partition as defined in \eqref{h}. Then there exists a constant $\hat{C} > 0$, independent of $h$, such that
    \begin{equation}
        \max_{1 \leq k \leq K} d_{N}(\mathcal{M}_{\mathcal{N},k}) \leq \hat C h^{N^{1/d}}.
    \label{hat C}
    \end{equation}
\end{proposition}

\begin{proof}
    We take $\rho_h=(\rho_{h,j})_{j=1}^d$ such that
    \begin{equation}
        \rho_{h,j}=\zeta_h :=1+\frac{2\varepsilon}{h}>1, \quad 1\leq j \leq d.
        \label{rho}
    \end{equation}
    It is clear that this choice of $\rho_h$ ensures that the polydisc $\mathcal{P}_{\rho_h}$ is contained in the open set $\hat{\mathcal{O}}$ for all subdomains $\mathcal{D}_k$. Let $\{\nu^{(n)}\}_{n \in \mathbb{N}}$ denote the multi-indices $\nu$ such that $\rho_h^{-\nu}=\zeta_h^{-|\nu|}$ are arranged in a non-increasing order. Let $\Lambda_N$ be a set containing the indices $\nu$ corresponding to the $N$ largest values of $\zeta_h^{-|\nu|}$. We consider the specific threshold
    \begin{equation}
        k_h= \hat{N} \lambda_{h}, 
        \label{kh}
    \end{equation}
    where
    \begin{equation}
        \hat{N}:=\left|\nu^{(N)}\right| +1, \quad \lambda_h:= \text{ln} (\zeta_h).
        \label{lambda def}
    \end{equation}
    Then, we have 
    \begin{equation}
        \Lambda_N\subseteq\left\{\nu \in \mathbb{N}^d: \zeta_h^{-|\nu|} > e^{-k_h}\right\}=:S_{k_h}.
    \end{equation}
    Equivalently, 
    \begin{equation}
        S_{k_h}:=\left\{\nu \in \mathbb{N}^d: \left|\nu\right| \lambda_h < k_h\right\}.
        \label{k}
    \end{equation}
    Due to ties in the values of $\zeta_h^{-|\nu|}$, the set $S_{k_h}$ may contain more than $N$ indices. In this case, the set $\Lambda_N$ s obtained by discarding indices corresponding to the smallest values of $\zeta_h^{-|\nu|}$ in $S_{k_h}$ so that $|\Lambda_N| = N$.
    
    The cardinality of $\Lambda_{N}$ is bounded by the number of all partial derivatives up to order $\hat{N}-1$ which is the maximum number of derivatives that could appear in the Taylor expansion if we consider the $N$ largest values of $\zeta_h^{-|\nu|}$. Then,
    \begin{align}
        N&\leq \sum_{n=0}^{\hat{N}-1} {n+d-1 \choose d-1} = \frac{\left(\hat{N}-1+d\right)!}{\left(\hat{N}-1\right)! \ d!}\\
        &= \frac{\left(\hat{N}-1+d\right) \left(\hat{N}-1+(d-1)\right) \cdots \left(\hat{N}\right)}{d!}\\
        &=\left(\frac{\hat{N}-1}{d} +1\right)\left(\frac{\hat{N}-1}{d-1} +1 \right) \cdots (\hat{N})\\
        &\leq (\hat{N})^d\\
        &=\left( \frac{k_h}{\lambda_h}\right)^d
        \label{est N}
    \end{align}
    
    In general, sets $\Lambda_n$ of arbitrary size $n$ corresponding to a threshold $k$ consist of at most all integer lattice points inside the simplex bounded by the coordinate hyperplanes together with the hyperplane $\lambda_{h} \sum_{j=1}^d t_j = k$. These sets are downward closed and their cardinality is bounded by the volume of the following continuous simplex
    \begin{equation}
        T_{k}:=\left\{\left(t_1, \ldots, t_d\right) \in \mathbb{R}^d: t_j \geq-1, j=1, \ldots, d, \text { and } \sum_{j=1}^d t_j \leq \frac{k}{\lambda_{h}} \right\}.
    \end{equation}
    Then,
    \begin{align}
        |\Lambda_n|=|S_{k}|\leq\left|T_k\right|&=\frac{1}{d!} \left(\frac{k+d \lambda_h}{\lambda_h}\right)^d.
        \label{set size}
    \end{align}
    
    For any $\mu \in \mathcal{D}$ there exists a subdomain $\mathcal{D}_k \subset \mathcal{D}$ such that $\mu \in \mathcal{D}_k$. Therefore, the approximation error when keeping only the $N$ terms with indices in $\Lambda_N$ is given by 
    \begin{align}
          \left\|u_{\mathcal{N}}(\mu)-\sum_{\nu \in \Lambda_N} t_{k,\nu} \mu_k^\nu \right\|_V 
        & \leq \sum_{\nu \notin S_{k_h}} \left\|t_{k,\nu} \right\|_V \\
        & \leq C\sum_{\nu \notin S_{k_h}} \zeta_h^{-|\nu|} \\
        & \leq C \sum_{l \geq k_h } e^{-l} \left|\left\{\nu: e^{-l-1} < \zeta_h^{-|\nu|} \leq e^{-l}\right\}\right|\\ %because we can only deal with (e^{-l-1} < \zeta_h^{-|\nu|} ) through S_k
        & \leq C \sum_{l \geq k_h } e^{-l} \left|S_{l+1}\right|.
    \end{align}
    Using the estimate \eqref{set size}, we obtain
    \begin{equation}
          \left\|u_{\mathcal{N}}(\mu)-\sum_{\nu \in \Lambda_N} t_{k,\nu} \mu_k^\nu \right\|_V  \leq \frac{C}{d! \ (\lambda_h)^d} \sum_{l \geq k_h } e^{-l} \left(l+1+ d\lambda_h\right)^d.
        \label{err_1}
    \end{equation}
    From \eqref{kh}, estimate \eqref{err_1} becomes
    \begin{align}
          \left\|u_{\mathcal{N}}(\mu)-\sum_{\nu \in \Lambda_N} t_{k,\nu} \mu_k^\nu \right\|_V &  \leq \frac{C \ \hat{N}^d}{d! \  (k_h)^d} \sum_{l \geq k_h } e^{-l} \left(l+1+\hat{N}^{-1}d k_h\right)^d\\
        & \leq \frac{C \ \hat{N}^d}{d!} \frac{\Phi\left(e^{-1},-d,(\hat{N}^{-1}d+1)k_h+1\right)}{(k_h)^d} e^{-k_h},
    \end{align}
    where $\Phi$ is the Hurwitz Lerch transcendent which converges for these values and gives a polynomial of degree $d$ in $k_h$. Consequently, the ratio $\frac{\Phi}{(k_h)^d}$ is uniformly bounded as $h\rightarrow 0$ $\left(k_h \rightarrow \infty\right)$. As a result, there exisits a constant $C_1$ such that
    \begin{align}
          \left\|u_{\mathcal{N}}(\mu)-\sum_{\nu \in \Lambda_N} t_{k,\nu} \mu_k^\nu \right\|_V &\leq C_1 e^{-k_h}\\
        &\leq C_1 e^{-\lambda_hN^{1/d}},
    \end{align}
    where $C_1$ is a constant that depends on $d$ and $N$, and the second inequality comes from estimate \eqref{est N}. From the definitions \eqref{rho} and \eqref{lambda def}, it can be inferred that 
    \begin{align}
        e^{-\lambda_h}&=\frac{1}{\zeta_h}= \frac{h}{h +2\varepsilon}\leq \frac{h}{2\varepsilon}.
    \end{align} 
    Therefore, we obtain the error estimate
    \begin{equation}
          \max_{1 \leq k \leq K}\sup_{\mu \in \mathcal{D}_k} \left\|u_{\mathcal{N}}(\mu)-\sum_{\nu \in \Lambda_N} t_{k,\nu} \mu_k^\nu \right\|_V \leq \hat C h^{N^{1/d}}.
    \end{equation}
    As a result, the Kolmogorov $N$-width inherits the same convergence rate which completes the proof
    
\end{proof}

\end{document}
